\documentclass[french,11pt]{book}
\usepackage{teach}
\usepackage[french]{babel}
\usepackage{amsmath,amssymb}
\usepackage{array}
\newcommand{\R}{\mathbb{R}}
\begin{document}
\begin{center}
{\Large \textbf{Devoir surveillé -- Variables aléatoires}}\\[2mm]
Terminale technologique
\end{center}

\section*{Devoir surveillé}
\begin{enumerate}
  \item Une loterie propose les gains 0 euro, 4 euros et 20 euros avec les probabilités 0,7 ; 0,25 ; 0,05. Calculer le gain moyen.
  \item Le billet coûte 2 euros. Déterminer l'espérance du gain algébrique et dire si le jeu est favorable au joueur.
  \item Calculer l'écart type de la variable de gain brut et interpréter la dispersion.
\end{enumerate}

\end{document}
